\section{Summary of Work}
\label{sec:conclusion:summary}

This thesis introduced a modular, rule-based railway interlocking system aimed at modernizing the control of train movements in station environments. Rooted in the limitations of traditional mechanical and relay-based systems, the work explored a software-first approach to safety logic, flexibility, and scalability.

The system was designed using Qt C++ and PostgreSQL, enabling real-time route validation, dynamic signal aspect control, and simulation-based testing. Key architectural components include a rule engine for logical enforcement, modular device managers for signals, point machines, and track circuits, and a simulation interface for operator interaction.

Through rigorous simulation, test cases, and performance benchmarks, the system demonstrated its ability to perform safe and deterministic decisions under both nominal and fault-induced conditions. The evaluation confirmed its potential as a foundation for future deployments in small to medium-sized railway stations.
